\section{Determining the monomial and the preferred framing}\label{pl:determining-the-monomial-and-the-preferred-framing}

By Section~\ref{pl:the-constant-field-of-the-source-equations} the discrepancy is a rational function of \(M\).
By Section~\ref{pl:regular-reconstruction-at-every-nonzero-main-value} its pullback is a regular unit along a local section
through every \(M_0\ne0\). If it had a zero or pole at \(M_0\),
that nonzero order would persist on the section. Therefore it has
no zero or pole on \(\mathbb C^*\), and

\begin{equation}\label{pl:F1}
\frac{L_{\mathrm{aug}}^2}{\prod_e f_e^2}=cM^k,\qquad
                  c\ne0,\quad k\in\mathbb Z.
\end{equation}

This is the elementary classification of units of
\(\mathbb C[M,M^{-1}]\). At the complete main cusp every
\(f_e=1\) and \(L_{\mathrm{aug}}^2=1\), so \(c=1\).

The calculation in Section~\ref{pl:cusp} is a marked face-word calculation on the same original exterior:

\begin{equation}\label{pl:F2}
\tau_{\mathrm{main}}
      =a_{\mathrm{aug}}+2i\sum_e\cot(\pi x_e),\qquad
 X_e=e^{2\pi ix_e},\qquad
 a_{\mathrm{aug}}=-w_\epsilon-\sum_v\epsilon_v.
\end{equation}

The integer includes the retained half twists and the preferred
zero-linking correction. It is not inferred from \eqref{pl:F1}.

For completeness, the cusp shape equals the logarithmic slope
of this geometric eigenvalue curve. Near a noncentral parabolic
meridian write a commuting determinant-one longitude as
\(aI+b\mu\). Its eigenvalues are \(a+bM^{\pm1}\).
At \(M=1\), its projective translation divided by the
meridian translation is \(b/(a+b)\). Differentiating the
determinant-one relation shows that
\(d\log L/d\log M\) has the same value. Thus \eqref{pl:F2}
computes this derivative with the specified marking.

Since \(f_e(1,X)=1\) for every \(X\), its first
derivative receives no contribution from the derivative
of the wire:

\begin{equation}\label{pl:F3}
\left.\frac{d\log f_e}{d\log M}\right|_{M=1}
       =2\frac{1+X_e}{1-X_e}
       =2i\cot(\pi x_e).
\end{equation}

Equations \eqref{pl:F1}--\eqref{pl:F3} give \(k=2a_{\mathrm{aug}}\), proving
\eqref{pl:T4}. Tangent agreement fixes the exponent only after the
preceding argument has excluded every other rational factor.

Now let \(\ell_v=\operatorname{lk}(K,C_v)\).
The crossing disk has exactly two main punctures, so
\(\ell_v=0\) for opposite local strand directions and
\(\ell_v=\pm2\) for parallel directions. In the augmented
exterior,

\begin{equation}\label{pl:F4}
[\lambda_{\mathrm{aug}}]=\sum_v\ell_v[\mu_v],
             \qquad [\lambda_v]=\ell_v[\mu_K].
\end{equation}

After filling \(\mu_v+p_v\lambda_v\), its meridian has
class \(-p_v\ell_v[\mu_K]\). The original main-boundary
word that becomes the preferred filled longitude is therefore

\begin{equation}\label{pl:F5}
\lambda_{\mathrm{fill}}
       =\lambda_{\mathrm{aug}}\,
                       \mu_K^{\sum_vp_v\ell_v^2}.
\end{equation}

This is a primitive integral basis change on the same main
torus. It is not an assertion that the augmented longitude
was already zero linking after surgery.



\begin{figure}[tb]
\centering\begin{tikzpicture}[>=Stealth,scale=.82,font=\small]
\coordinate (O) at (0,0);\coordinate (U) at (.6,0);
\coordinate (A) at (.65,2.3);\coordinate (F) at (3.05,2.3);
\fill[blue!8] (O)--(U)--($(U)+(A)$)--(A)--cycle;
\fill[orange!12] (O)--(U)--($(U)+(F)$)--(F)--cycle;
\draw[blue!65!black] (O)--(U)--($(U)+(A)$)--(A)--cycle;
\draw[orange!75!black] (O)--(U)--($(U)+(F)$)--(F)--cycle;
\draw[->,thick] (O)--(U) node[midway,below] {$\mu_K$};
\draw[->,thick,blue!65!black] (O)--(A) node[pos=.7,left] {$\lambda_{\rm aug}$};
\draw[->,thick,orange!75!black] (O)--(F) node[pos=.58,right] {$\lambda_{\rm fill}$};
\draw[->,dashed] (A)--(F) node[midway,above] {$4\mu_K$};
\node[align=left,anchor=west] at (5,.95) {Same cusp and meridian.\\Different preferred basis\\after disk filling.};
\end{tikzpicture}

\caption{The image of the augmented longitude and the preferred filled longitude differ by the integral meridian shear in \eqref{pl:F5}. The example shows $p=1$, $\ell=2$, hence shear four; in general it is $\sum_vp_v\ell_v^2$. Both parallelograms generate the same filled cusp lattice.}
\end{figure}

Write \(o_v=1\) for parallel and \(o_v=-1\) for
opposite strand directions in the fixed physical convention.
Then
\[
 \ell_v^2=2(1+o_v),\qquad
 w_\epsilon=\sum_vo_v\epsilon_v,\qquad
 w(D)=w_\epsilon-2\sum_vo_vp_v.
\]
Consequently

\begin{equation}\label{pl:F6}
\begin{aligned}
 a_{\mathrm{aug}}+\sum_vp_v\ell_v^2
  &=-w_\epsilon-\sum_v\epsilon_v+2\sum_vp_v(1+o_v)\\
  &=-w(D)-\sum_v(\epsilon_v-2p_v)=a(D).
\end{aligned}
\end{equation}

Combining \eqref{pl:T4} with \eqref{pl:F5}--\eqref{pl:F6} proves the squared filled
identity \eqref{pl:T6}, with the original primitive slopes and
preferred longitude.
